\documentclass[10pt, t, aspectratio = 1610]{beamer}
\usepackage{./template/beamertheme}
\usepackage{amsmath}
  \newcommand{\e}{\mathrm{e}}
\usefonttheme[onlymath]{serif}
\usepackage[UTF8, linespread = 1.5]{ctex}
\usepackage{parskip}
\usepackage[many, minted]{tcolorbox}
\makeatletter
\renewenvironment{quote}{%
    \begin{tcolorbox}%
    [
        top = 3pt,
        bottom = 3pt,
        left = 3pt,
        right = 3pt,
        colback = beamer@barcolor!10!white,
        colframe = beamer@headercolor,
        colupper = beamer@normaltextcolor,
    ]%
}%
{\end{tcolorbox}}

\def\inlinebox@true{true}
\define@key{inlinebox}{frame rule color} {\def\inlinebox@framerulecolor{#1}}
\define@key{inlinebox}{frame back color} {\def\inlinebox@framebackcolor{#1}}
\define@key{inlinebox}{frame text color} {\def\inlinebox@frametextcolor{#1}}
\define@key{inlinebox}{frame rule width} {\def\inlinebox@framerulewidth{#1}}
\define@key{inlinebox}{banner width}     {\def\inlinebox@bannerwidth{#1}}
\define@key{inlinebox}{show banner}[true]{\def\inlinebox@showbanner{#1}}
\define@key{inlinebox}{banner text color}{\def\inlinebox@bannertextcolor{#1}}
\define@key{inlinebox}{banner back color}{\def\inlinebox@bannerbackcolor{#1}}
\define@key{inlinebox}{banner text}      {\def\inlinebox@bannertext{#1}}
\NewDocumentCommand{\inlinebox}{O{} m}{%
  \setkeys{inlinebox}{%
    frame rule color  = black,
    frame back color  = white,
    frame text color  = black,
    frame rule width  = 0.4pt,
    banner width      = 8pt,
    show banner       = false,
    banner text color = white,
    banner back color = black,
    banner text       = BAN,
    #1
  }%
  \tcbox[%
    enhanced,
    tcbox raise base,
    nobeforeafter,
    boxrule           = \inlinebox@framerulewidth,
    top               = -1pt,
    bottom            = -1pt,
    right             = -1pt,
    arc               = 1pt,
    left              = \ifx\inlinebox@showbanner\inlinebox@true\inlinebox@bannerwidth-2pt\else-1pt\fi,
    colframe          = \inlinebox@framerulecolor,
    coltext           = \inlinebox@frametextcolor,
    colback           = \inlinebox@framebackcolor,
    before upper      = {\vphantom{蛤dg}},
    overlay           = {%
      \begin{tcbclipinterior}
        \ifx\inlinebox@showbanner\inlinebox@true
          \fill[\inlinebox@bannerbackcolor] (frame.south west) rectangle node[text = \inlinebox@bannertextcolor, scale = 0.4, font = \sffamily\bfseries, rotate = 90] {\inlinebox@bannertext} ([xshift = \inlinebox@bannerwidth]frame.north west);
        \fi
      \end{tcbclipinterior}%
    }%
  ]{#2}%
}

\newtcblisting{pythoncode}{
    listing engine = minted,
    minted style = colorful,
    minted language = python,
    minted options = {%
        fontsize = \fontsize{9}{9}\selectfont,
        linenos,
        numbersep = 3mm,
        autogobble,
        breaklines,
        breakafter = d},
    colback = beamer@barcolor!5!white,
    colframe = beamer@headercolor,
    listing only,
    left = 5mm,
    top = 0pt,
    bottom = 0pt,
    enhanced,
    overlay = {%
        \begin{tcbclipinterior}\fill[beamer@headercolor!20!white] (frame.south west)
            rectangle ([xshift=5mm]frame.north west);\end{tcbclipinterior}}%
}
\newcommand{\inlinepython}[1]{\inlinebox[frame rule color = black!30!white]{\mintinline{python}{#1}}}
\newcommand{\inlineframe}[1]{\inlinebox[frame rule color = black]{#1}}
\makeatother



\usepackage{enumitem}
  \setitemize[1]{itemsep = 0pt, partopsep = 0pt, parsep = 0pt, topsep = 0pt, leftmargin = 1em, label = $\bullet$}
  \setitemize[2]{itemsep = 0pt, partopsep = 0pt, parsep = 0pt, topsep = 0pt, leftmargin = 0pt, label = $\bullet$}
  \usepackage{tabu}
  \tabuDisableCommands{%
    \renewcommand\mintinline[2]{\texttt{\detokenize{#2}}}%
  }
  \tabulinesep = 3pt
\makeatletter
  \setlength{\parskip}{\medskipamount}

  \newcommand{\@minipagerestore}{\setlength{\parskip}{\medskipamount}}
  \makeatother
\title{TensorFlow 入门}
\subtitle{搭建自己的神经网络}
\date{\today}
\author{%
    Qiqi%
}

\institute%
{%
    2333 $\cdot$ 2333\\%
    23333%
}

\pgfdeclareimage[height = 1cm]{titlepagelogo}{./template/icon_in_cover.pdf}
\titlegraphic{%
  \pgfuseimage{titlepagelogo}%
}

\newcommand{\insertcover}{{
    \aauwavesbg%
    \begin{frame}[plain, noframenumbering]
         \titlepage
    \end{frame}
}}

\newcommand{\insertthanks}{{
    \aauwavesbg
    \begin{frame}[plain, noframenumbering]
        \finalpage{\huge 谢谢, 欢迎提问\vspace{-10pt}}
    \end{frame}
}}

\usepackage{pgfplots}
\usetikzlibrary{arrows.meta}
\tikzset
{
  > = {Stealth[length=5pt]},
}

\begin{document}
\insertcover

\begin{frame}{目录}
    \tableofcontents
\end{frame}

%!TEX root = ../main.tex


\section{TensorFlow 简介}

\begin{frame}{\secname}
\begin{center}
  \vfill
  \raisebox{-7pt}{\includegraphics[width = 8cm]{./template/icon_in_cover.pdf}} \scalebox{3.2}{\bfseries 简介}\vspace{10pt}

  \begin{varwidth}{\textwidth}
  \large
  \tableofcontents[sectionstyle = hide/hide, subsectionstyle = show/show/hide]
  \end{varwidth}
  \vfill
\end{center}
\end{frame}

\subsection{TensorFlow 是什么}
\begin{frame}{\insertsection}{\insertsubsection}
    TensorFlow 是一款通过数据流图进行数值计算的开源库. TensorFlow 最早是被 Google Brain 的研究者和工程师开发出来的用于机器学习和深度神经网络的研究. 但是这个系统也同样适用于其他领域.

    TensorFlow 具有很多非常好的特性:%
    %
    \begin{itemize}
        \item \textbf{易用性}, TensorFlow 封装了很多复杂运算, 使用者只需构建计算图.
        \item \textbf{可移植性}, TensorFlow 支持多种平台, 系统以及硬件.
        \item \textbf{自动求导}, TensorFlow 支持自动求导, 给采用基于梯度下降的学习方法带来了极大便利.
        \item \textbf{支持多种语言}, TensorFlow 提供方便易用的 Python 接口, 也支持 C++, Java, Go 等语言.
        \item \textbf{性能最大化}, 可以部署到不同硬件, 并使用线程, 队列, 异步计算来最大化硬件使用效率.
    \end{itemize}
\end{frame}

\subsection{TensorFlow 的主要技术特性}
\begin{frame}{\insertsection}{\insertsubsection}
    \begin{table}[!htb]
        \centering
        \begin{tabu}{X[l, -1]|X[l]}
            \tabucline[1pt]{-}
            编程模型 & Dataflow-Like Model\\\hline
            语言     & Python, C++, Go, Rust, Haskell, Java, Julia, JavaScript, R\\\hline
            部署     & Code-once, run everywhere\\\hline
            计算资源 & CPU, GPU, TPU\\\hline
            实现方式 & Local Implementation, Distributed Implementation\\\hline
            平台支持 & Google Cloud Platform, Hadoop File System\\\hline
            数学表达 & Math Graph Expression, Auto Differentiation\\\hline
            优化     & Common Subexpression Elimination,\\
                     & Asynchronous Kernel Optimization,\\
                     & Communication Optimization,\\
                     & Model Parallelism, Data Parallelism, Pipeline\\
            \tabucline[1pt]{-}
        \end{tabu}
    \end{table}
\end{frame}


% \begin{frame}{\insertsection}{TensorFlow 的工作原理}
%     \begin{minipage}[m]{0.5\textwidth}
%     TensorFlow 是用数据流图对计算过程进行描述的. 在数据流图中, 节点通常代表数学运算, 边表示节点之间的某种联系, 边负责在节点之间传输多维数据, 即 Tensor.

%     节点可以被分配到多个计算设备上, 可以异步和并行的进行操作. 因为是有向图, 所以只能等待之前的节点运行结束, 当前节点才能执行操作. 如图所示是一个简单的计算流图, 且它的表达式为%
%     \[
%         \bigg(2.2 - \Big(\frac{x}{11}\Big)\bigg) + \big(7 \times \cos(y)\big)\text{.}
%     \]
%     \end{minipage}\hfill
%     \begin{minipage}[m]{0.4\textwidth}
%     \begin{figure}[!htb]
%         \centering
%         \scalebox{0.8}{%!TEX root = ../main.tex

\begin{tikzpicture}[%
    node/.style = {circle, draw, minimum size = 0.8cm, align = flush center, inner sep = 0pt, fill = black, text = white},
    y = 0.8cm]

    \node[node] (x) at (5, 0) {$x$};
    \node[node] (11) at (6, 0) {$11$};
    \node[node] (y) at (9, 0) {$y$};
    \node[node] (div) at (5.5, 2) {$\div$};
    \node[node] (22) at (3.2, 2) {$2.2$};
    \node[node] (7) at (7.5, 2) {$7$};
    \node[node] (cos) at (9.5, 2) {$\cos$};
    \node[node] (sub) at (4.5, 4) {$-$};
    \node[node] (times) at (8.5, 4) {$\times$};
    \node[node] (plus) at (6.5, 6) {$+$};

    \draw[->] (x) -- (div);
    \draw[->] (11) -- (div);
    \draw[->] (y) -- (cos);
    \draw[->] (22) -- (sub);
    \draw[->] (div) -- (sub);
    \draw[->] (7) -- (times);
    \draw[->] (cos) -- (times);
    \draw[->] (sub) -- (plus);
    \draw[->] (times) -- (plus);
\end{tikzpicture}
}
%     \end{figure}
%     \end{minipage}
% \end{frame}

%!TEX root = ../main.tex

\section{TensorFlow 基本概念}
\subsection{张量 Tensor}
\begin{frame}{\insertsection}{\insertsubsection}
张量最初是个物理学的概念,
\begin{quote}
    A tensor is something that transforms like a tensor\\
    一个量, 在不同的参考系下按照某种特定的法则进行变换, 就是张量.\\
    \rule{0pt}{0pt}\hfill  --- A. Zee, "Einstein Gravity in a Nutshell"
\end{quote}

在 TensorFlow 中, 可以将张量理解成多维数组, 张量穿梭在数据流图中, 当经过节点时, 会被节点按照特定法则变换成另一个张量, 这个就是 TensorFlow 名字的由来.

张量的维数被称为张量的阶, 我们常见的标量, 其实是 $0$ 阶张量, 矢量就是 $1$ 阶张量, 矩阵就是 $2$ 阶张量, 一幅图片, 在 TensorFlow 中可以是 $3$ 阶张量.
\end{frame}

\begin{frame}{\insertsection}{\insertsubsection}
TensorFlow 文档中使用了三种记号来方便地描述张量的维度: 阶, 形状以及维数. 下表展示了他们之间的关系:\vspace{10pt}

\begin{table}
  \centering
  \begin{tabu}{clcl}
  \tabucline[1pt]{-}
  \rowfont{\bfseries}
  阶       & 形状                                     & 维数     & python 代码\\
  \hline
  $0$      & \texttt{[ ]}                             & $0$      & \inlinepython{t0 = 483}\\
  $1$      & \texttt{[$D_0$]}                         & $1$      & \inlinepython{t1 = [1.1, 2.2, 3.3]}\\
  $2$      & \texttt{[$D_0$, $D_1$]}                  & $2$      & \inlinepython{t2 = [[1, 2, 3], [4, 5, 6], [7, 8, 9]]}\\
  $3$      & \texttt{[$D_0$, $D_1$, $D_2$]}           & $3$      & \inlinepython{t3 = [[[1], [2]], [[3], [4]], [[5], [6]]]}\\
  $\vdots$ & $\vdots$                                 & $\vdots$ & $\vdots$\\
  $n$      & \texttt{[$D_0$, $D_1$, $\cdots$, $D_n$]} & $n$      & \inlinepython{tn = ...}\\
  \tabucline[1pt]{-}
  \end{tabu}
\end{table}
\end{frame}

\begin{frame}{\insertsection}{\insertsubsection}
除了维度, Tensors有一个数据类型属性. 可为一个张量指定下列数据类型的任意一个类型:\vspace{10pt}

\begin{table}
  \centering
  \begin{tabu}{ll|ll}
  \tabucline[1pt]{-}
  \rowfont{\bfseries}
  类型 & 描述 & 类型 & 描述 \\
  \hline
  \inlinepython{tf.float32} & 32 位浮点数     & \inlinepython{tf.string}    & 可变长度的字节数组\\
  \inlinepython{tf.float64} & 64 位浮点数     & \inlinepython{tf.bool}      & 布尔型\\
  \inlinepython{tf.int64}   & 64 位有符号整型 & \inlinepython{tf.complex64} & 由两个 32 位浮点数组成的复数\\
  \inlinepython{tf.int32}   & 32 位有符号整型 & \inlinepython{tf.qint32}    & 用于量化操作的 32 位有符号整型\\
  \inlinepython{tf.int16}   & 16 位有符号整型 & \inlinepython{tf.qint8}     & 用于量化操作的 8 位有符号整型\\
  \inlinepython{tf.int8}    & 8 位有符号整型  & \inlinepython{tf.quint8}    & 用于量化操作的 8 位无符号整型\\
  \inlinepython{tf.uint8}   & 8 位无符号整型  & & \\
  \tabucline[1pt]{-}
  \end{tabu}
\end{table}
\end{frame}

\begin{frame}[fragile]{\insertsection}{\insertsubsection}
TensorFlow 中张量包括 \inlinepython{constant}, \inlinepython{placeholder} 和 \inlinepython{Variable}.

\begin{onlyenv}<2>
\inlinepython{tf.constant()} 函数提供在 TensorFlow 中定义不可更改张量的方法, 定义如下所示.

\begin{pythoncode}
def constant(value, dtype=None, shape=None, name="Const", verify_shape=False)
\end{pythoncode}

\begin{itemize}
\item \inlinepython{value}, 符合 TensorFlow 中定义的数据类型的常数值或者常数列表;
\item \inlinepython{dtype}, 数据类型, 可选;
\item \inlinepython{shape}, 常量的形状, 可选;
\item \inlinepython{name}, 常量的名字, 可选;
\item \inlinepython{verify_shape}, 常量的形状是否可以被更改, 默认不可更改.
\end{itemize}

除了直接赋值以外, 还可使用 \inlinepython{tf.ones()}, \inlinepython{tf.zeros()} 等初始化张量的方法.
\end{onlyenv}

\begin{onlyenv}<3>
\inlinepython{tf.placeholder()} 函数提供在 TensorFlow 中定义占位张量的方法, 定义如下所示.

\begin{pythoncode}
def placeholder(dtype, shape=None, name=None)
\end{pythoncode}

\begin{itemize}
\item \inlinepython{dtype} 指定占位张量的数据类型, 可以是 TensorFlow 中的数据类型, 如常用的 \inlinepython{tf.float32}, \inlinepython{tf.float64} 等数值类型;
\item \inlinepython{shape} 表示数据类型, 默认的 \inlinepython{None} 是一个一维的数值, \inlinepython{shape = [None, 5]}, 表示行不定, 列是 \inlinepython{5};
\item \inlinepython{name} 是张量名称;
\end{itemize}

占位变量是一种 TensorFlow 用来解决读取大量训练数据问题的机制, 它允许你现在不用给它赋值, 随着训练的开始, 再把训练数据传送给训练网络学习.
\end{onlyenv}

\begin{onlyenv}<4->
TensorFlow 中变量是通过 \inlinepython{Variable} 类来实现的, 初始化函数 \inlinepython{tf.Variable()} 定义如下.

\begin{pythoncode}
def __init__(self, initial_value=None, trainable=True, collections=None, validate_shape=True, caching_device=None, name=None, variable_def=None, dtype=None, expected_shape=None, import_scope=None)
\end{pythoncode}


\only<4>{
\begin{itemize}
\item \inlinepython{initial_value}, 初始值, 必填, 张量或可以转换为张量的 Python 对象;
\item \inlinepython{trainable}, 如果参数 \inlinepython{trainable} 的值为 \inlinepython{True}, 则默认值也将变量添加到图形中集合 \inlinepython{GraphKeys.TRAINABLE_VARIABLES}。这个集合用作 \inlinepython{Optimizer} 类使用的默认变量列表;
\item \inlinepython{collections}, 新的变量被添加到这些集合。默认为 \inlinepython{[GraphKeys.GLOBAL_VARIABLES]};
\end{itemize}
}
\only<5>{
\begin{itemize}
\item \inlinepython{validate_shape}, 如果 \inlinepython{False}, 允许变量用初始化未知形状的值。如果 \inlinepython{True}, 默认的形状 \inlinepython{initial_value} 必须是已知的;
\item \inlinepython{caching_device}, 指定变量的默认缓存位置;
\item \inlinepython{name}, 是张量名称;
\end{itemize}
}
\only<6>{
\begin{itemize}
\item \inlinepython{variable_def}, 指定 \inlinepython{VariableDef} 的协议缓冲区;
\item \inlinepython{dtype}, 指定张量的数据类型;
\item \inlinepython{expected_shape}, 是张量的形状, 如果设置, \inlinepython{initial_value} 需要符合这个形状;
\item \inlinepython{import_scope}, 张量名称前缀.
\end{itemize}
}

\end{onlyenv}
\end{frame}

\subsection{节点 Operator}
\begin{frame}{\insertsection}{\insertsubsection}
在 TensorFlow 中, 每一个节点代表着一个操作, 一般用来表示施加的的数学运算, 也可以表示数据输入的起点以及输出的终点. 下面是一些重要的操作:

\begin{table}
  \centering
  \begin{tabu}{ll|ll}
  \tabucline[1pt]{-}
  \rowfont{\bfseries}
   操作 & 描述 &  操作 & 描述 \\
  \hline
  \inlinepython{tf.add(x, y, name=None)} & 求和 & \inlinepython{tf.sign(x, name=None)}   & 返回符号\\
  \inlinepython{tf.sub(x, y, name=None)} & 减法 & \inlinepython{tf.neg(x, name=None)}    & 取负 ($y = -x$)\\
  \inlinepython{tf.mul(x, y, name=None)} & 乘法 & \inlinepython{tf.square(x, name=None)} & 计算平方 ($y = x^2$)\\
  \inlinepython{tf.div(x, y, name=None)} & 除法 & \inlinepython{tf.round(x, name=None)}  & 求最接近的整数\\
  \inlinepython{tf.mod(x, y, name=None)} & 取模 & \inlinepython{tf.sqrt(x, name=None)}   & 开根号 ($y = \sqrt{x}$)\\
  \inlinepython{tf.pow(x, y, name=None)} & 乘幂 & \inlinepython{tf.abs(x, name=None)}    & 绝对值\\
  \inlinepython{tf.inv(x, name=None)}    & 取反 & \inlinepython{tf.exp(x, name=None)}    & 计算 $\e$ 的次方\\
  \tabucline[1pt]{-}
  \end{tabu}
\end{table}
\end{frame}

\subsection{数据流图 Graph}
\begin{frame}{\insertsection}{\insertsubsection}
\vspace{10pt}
\begin{minipage}[m]{0.5\textwidth}
    TensorFlow 是用数据流图对计算过程进行描述的. 在数据流图中, 节点代表数学运算, 边表示节点之间的某种联系, 负责在节点之间传输即张量.

    节点可以被分配到多个计算设备上, 可以异步和并行的进行操作. 因为是有向图, 所以只能等待之前的节点运行结束, 当前节点才能执行操作.

    如图所示是一个简单的计算流图.%
    \end{minipage}\hfill
    \begin{minipage}[m]{0.4\textwidth}
    \inlineframe{\scalebox{0.8}{%!TEX root = ../main.tex

\begin{tikzpicture}[%
    node/.style = {circle, draw, minimum size = 0.8cm, align = flush center, inner sep = 0pt, fill = black, text = white},
    y = 0.8cm]

    \node[node] (x) at (5, 0) {$x$};
    \node[node] (11) at (6, 0) {$11$};
    \node[node] (y) at (9, 0) {$y$};
    \node[node] (div) at (5.5, 2) {$\div$};
    \node[node] (22) at (3.2, 2) {$2.2$};
    \node[node] (7) at (7.5, 2) {$7$};
    \node[node] (cos) at (9.5, 2) {$\cos$};
    \node[node] (sub) at (4.5, 4) {$-$};
    \node[node] (times) at (8.5, 4) {$\times$};
    \node[node] (plus) at (6.5, 6) {$+$};

    \draw[->] (x) -- (div);
    \draw[->] (11) -- (div);
    \draw[->] (y) -- (cos);
    \draw[->] (22) -- (sub);
    \draw[->] (div) -- (sub);
    \draw[->] (7) -- (times);
    \draw[->] (cos) -- (times);
    \draw[->] (sub) -- (plus);
    \draw[->] (times) -- (plus);
\end{tikzpicture}
}}
    \[
        \text{表达式: }\bigg(2.2 - \Big(\frac{x}{11}\Big)\bigg) + \big(7 \times \cos(y)\big)\text{.}
    \]
\end{minipage}
\end{frame}

\begin{frame}[fragile]{\insertsection}{\insertsubsection}
\begin{pythoncode}
import tensorflow as tf
import numpy as np

x1 = tf.placeholder(tf.float32, [None, 1])
w1 = tf.Variable(tf.random_normal([1, 10]))
b1 = tf.Variable(tf.ones([1, 10]))

x2 = tf.nn.relu(tf.matmul(x1, w1) + b1)
w2 = tf.Variable(tf.random_normal([10, 1]))
b2 = tf.Variable(tf.ones([1, 1]))

predict_values = tf.matmul(x2, w2) + b2
actual_values = tf.placeholder(tf.float32, [None, 1])
loss = tf.reduce_mean(tf.reduce_sum(tf.square(actual_values - predict_values), reduction_indices = [1]))
\end{pythoncode}

\vspace{-7.5cm}\hspace{10.5cm}\scalebox{0.53}{\begin{tikzpicture}[%
  tensor/.style = {rectangle, inner sep = 2pt, rounded corners = 0.4cm, draw, align = flush center, minimum height = 0.8cm, minimum width = 0.8cm},
  operator/.style = {tensor, fill = black, draw = none, text = white, inner sep = 2pt},
  line width = 1pt,
  y = 1.1cm
]
\node[tensor, circle]   (x1) at (9, 0) {$x_1$};
\node[tensor, circle]   (w1) at (7, 0) {$W_1$};
\node[operator]         (t1) at (8, 1) {\mintinline{text}{tf.matmul}};
\node[tensor, circle]   (b1) at (6, 1) {$b_1$};
\node[operator]         (p1) at (7, 2) {\mintinline{text}{tf.add}};
\node[operator]         (r1) at (7, 3) {\mintinline{text}{tf.nn.relu}};
\node[tensor, circle]   (w2) at (5, 3) {$W_2$};
\node[operator]         (t2) at (6, 4) {\mintinline{text}{tf.matmul}};
\node[tensor, circle]   (b2) at (4, 4) {$b_2$};
\node[operator]         (p2) at (5, 5) {\mintinline{text}{tf.add}};
\node[tensor]           (pv) at (5, 6) {\mintinline{text}{prediction_values}};
\node[tensor]           (av) at (8, 6) {\mintinline{text}{actual_values}};
\node[operator]         (sb) at (6.5, 7.5) {\mintinline{text}{tf.sub}};
\node[operator]         (sq) at (6.5, 8.5) {\mintinline{text}{tf.square}};
\node[operator]         (rs) at (6.5, 9.5) {\mintinline{text}{tf.reduce_sum}};
\node[operator]         (rm) at (6.5, 10.5) {\mintinline{text}{tf.reduce_mean}};
\node[tensor]           (ls) at (6.5, 11.5) {\mintinline{text}{loss}};

\draw[->] (x1) -- (t1);
\draw[->] (w1) -- (t1);
\draw[->] (b1) -- (p1);
\draw[->] (t1) -- (p1);
\draw[->] (p1) -- (r1);
\draw[->] (w2) -- (t2);
\draw[->] (r1) -- (t2);
\draw[->] (b2) -- (p2);
\draw[->] (t2) -- (p2);
\draw[->] (p2) -- (pv);
\draw[->] (pv) -- (sb);
\draw[->] (av) -- (sb);
\draw[->] (sb) -- (sq);
\draw[->] (sq) -- (rs);
\draw[->] (rs) -- (rm);
\draw[->] (rm) -- (ls);
\end{tikzpicture}}
\end{frame}

\subsection{会话 Session}
\begin{frame}[fragile]{\insertsection}{\insertsubsection}
构造阶段完成后, 才能启动图。启动图的第一步是创建一个 Session 对象, 如果无任何创建参数, 会话构造器将启动默认图。会话会管理TensorFlow程序运行时的所有资源。当所有计算完成之后需要关闭会话来帮助系统回收资源，否则就可能出现资源泄露的问题。

\begin{minipage}[t]{0.48\textwidth}
\begin{pythoncode}
import tensorflow as tf

a = tf.constant([1.0, 2.0])
b = tf.constant([2.0, 3.0])
result = a + b

session = tf.Session()
print(session.run(result))
session.close()
\end{pythoncode}
\end{minipage}%
\hfill%
\begin{minipage}[t]{0.48\textwidth}
\begin{pythoncode}
import tensorflow as tf

a = tf.constant([1.0, 2.0])
b = tf.constant([2.0, 3.0])
result = a + b

with tf.Session() as session:
    print(session.run(result))
# session.close()
\end{pythoncode}
\end{minipage}
\end{frame}
%!TEX root = ../main.tex

\section{TensorFlow 实战}
\begin{frame}{\secname}
\vfill\hspace{20pt}
\begin{minipage}[m]{0.25\textwidth}
  \linespread{1}\large
  \tableofcontents[sectionstyle = hide/hide, subsectionstyle = show/show/hide]
\end{minipage}%
\hfill%
\begin{minipage}[m]{0.6\textwidth}
  \includegraphics[height = 7cm]{./icons/programming.pdf}
\end{minipage}%
\vfill
\end{frame}

\subsection{分类问题}
\begin{frame}{\insertsection}{\insertsubsection}
\begin{minipage}[m]{0.48\textwidth}
回到最开始的那个二分类问题, 看看用 \tensorflow{} 是如何解决这个问题的, $500$ 个点的分布如下图所示.

我们通过观察数据的特征, 仍然选用一样的神经网络, 如下图所示.

\centering
\scalebox{0.8}{%
\begin{tikzpicture}[%
  node/.style = {circle, draw, inner sep = 0pt, minimum size = 0.8cm},
  x = 1.4cm]
  \node[node] (x1) at (0, 2) {$x_1$};
  \node[node] (x2) at (0, 1) {$x_2$};
  \node[node] (b)  at (0, 0) {$1$};
  \node[node] (sum) at (2, 1) {$\sum$};
  \node[node] (step) at (4, 1) {$H(x)$};

  \draw[->] (x1) -- (sum) node[midway, fill = white, inner sep = 1pt, sloped] {$W_1$};
  \draw[->] (x2) -- (sum) node[midway, fill = white, inner sep = 1pt, sloped] {$W_2$};
  \draw[->] (b) -- (sum);
  \draw[->] (sum) -- (step);
  \draw[->] (step) -- ++(0:1.5cm) node[anchor = west] {$y$};
\end{tikzpicture}}
\end{minipage}%
\hfill%
\begin{minipage}[m]{0.48\textwidth}
\begin{tikzpicture}
  \begin{axis}[%
    xlabel = $x_1$,
    ylabel = $x_2$,
    height = 7.8cm,
    width = 7cm,
    title = $500$ 个点的分布图
  ]
    \addplot[only marks, discard if not={type}{1}, blue] table[col sep=comma, x = x1, y = x2]{./data/classification/training_data.dat};
    \addplot[only marks, discard if not={type}{0}, green] table[col sep=comma, x = x1, y = x2]{./data/classification/training_data.dat};
  \end{axis}
\end{tikzpicture}
\end{minipage}
\end{frame}

\begin{frame}[fragile]{\insertsection}{\insertsubsection}
这次我们不用计算神经网络的表达式, 直接照着神经网络写 \tensorflow{} 代码即可.

\begin{pythoncode}[fontsize = \fontsize{8}{8}\selectfont]
import tensorflow as tf
import numpy as np

# 定义 $\bm{x} = (x_1, x_2)$, $\bm{W} = (W_1, W_2)$ 和 $b = 1$
x = tf.placeholder(tf.float32, [None, 2])
w = tf.Variable(tf.zeros([2, 1]))
b = tf.constant([1.])

# 定义神经网络输出
predict_type = tf.sigmoid((tf.matmul(x, w) + b) * 10)
# 定义真实输出 $y = g(\bm{W}\bm{x} + 1)$
actual_type = tf.placeholder(tf.float32, [None, 1])
# 定义目标函数 $\ell(\bm{W}) =  -\sum_{i = 1}^{500} y_i\ln g(\bm{W}\bm{x}_i + 1) + (1 - y_i)\ln\big(1 - g(\bm{W}\bm{x}_i + 1)\big)$
loss = -tf.reduce_sum(actual_type * tf.log(predict_type) + (1 - actual_type) * tf.log(1 - predict_type))
\end{pythoncode}

\vspace{-6.3cm}\hfill%
\scalebox{0.8}{%
\begin{tikzpicture}[%
  node/.style = {circle, draw, inner sep = 0pt, minimum size = 0.8cm},
  x = 1.4cm, y = 1.4cm]
  \node[node] (x1) at (0, 2) {$x_1$};
  \node[node] (x2) at (0, 1) {$x_2$};
  \node[node] (b)  at (0, 0) {$1$};
  \node[node] (sum) at (2, 1) {$\sum$};
  \node[node] (step) at (4, 1) {$H(x)$};

  \draw[->] (x1) -- (sum) node[midway, fill = barcolor!5!white, inner sep = 1pt, sloped] {$W_1$};
  \draw[->] (x2) -- (sum) node[midway, fill = barcolor!5!white, inner sep = 1pt, sloped] {$W_2$};
  \draw[->] (b) -- (sum);
  \draw[->] (sum) -- (step);
  \draw[->] (step) -- ++(0:1.5cm) node[anchor = west] {$y$};
\end{tikzpicture}\ \ }
\end{frame}

\begin{frame}[fragile]{\insertsection}{\insertsubsection}
\begin{pythoncode}[fontsize = \fontsize{8}{8}\selectfont]
def load_training_data(path): # 加载训练数据 $\bm{x}$ 和 $\bm{y}$
  with open(path, 'r', encoding = 'utf8') as file:
    next(file)
    lines = [[float(item) for item in line.split()] for line in file]
    xs = np.array([line[:2] for line in lines])
    types = np.array([line[2] for line in lines])[:, None]
  return xs, types

with tf.Session() as session:
  session.run(tf.global_variables_initializer()) # 初始化变量
  train_step = tf.train.GradientDescentOptimizer(0.003).minimize(loss) # 梯度下降法
  xs, types = load_training_data('./training.dat') # 加载训练数据 $\bm{x}$ 和 $\bm{y}$
  for i in range(1000): # 训练 $1000$ 次
    session.run(train_step, feed_dict = {x:xs, actual_type:types}) # 填充训练数据
    # 打印目标函数值
    if i % 10 == 0: print(session.run(loss, feed_dict = {x:xs, actual_type:types}))
  print(session.run(w, feed_dict = feed_dict)) # 打印最终结果
\end{pythoncode}
\end{frame}

\begin{frame}[fragile]{\insertsection}{\insertsubsection}
最终, 参数 $W_1$ 和 $W_2$ 收敛于点 $(1.07, -3.10)$, 也是非常理想的结果.

我们还可以换一个目标函数, 用模型预测错误的数量作为目标函数, 其表达式如下所示:
\begin{flalign*}
&& \ell(\bm{W}) &= \sum_{\mathclap{x_1, x_2, y}} \big|H(W_1x_1+ W_2x_2 + 1)) - y\big|\text{,} &&\\
\text{由于绝对值不可微,}&& \text{上式} &= \sum_{\mathclap{x_1, x_2, y}} \big(H(W_1x_1+ W_2x_2 + 1)) - y\big)^2\text{,} &&\text{\hphantom{由于绝对值不可微,}}\\
\text{Sigmoid 函数替换,}&& &\approx\sum_{\mathclap{x_1, x_2, y}} \big(g(W_1x_1+ W_2x_2 + 1)) - y\big)^2\text{.} &&
\end{flalign*}

我们修改目标函数的表达式, 参数 $W_1$ 和 $W_2$ 收敛于点 $(1.07, -3.07)$, 效果拔群.
\begin{pythoncode}
loss = tf.reduce_sum(tf.square(actual_type - predict_type))
\end{pythoncode}
\end{frame}

\subsection{拟合问题}
\begin{frame}{\insertsection}{\insertsubsection}
\begin{minipage}[t]{0.41\textwidth}
有 300 个点, 其分布如下图所示, 现在的目标是用一个神经网络, 来拟合这 300 个点.

\begin{center}
\hspace{-10pt}%
\begin{tikzpicture}
\begin{axis}[width = 1.2\textwidth,
             height = 6.4cm]
\addplot[only marks] table[col sep = comma] {./data/fitting/training_data.dat};
\end{axis}
\end{tikzpicture}
\end{center}
\end{minipage}%
\hfill%
\begin{minipage}[t]{0.55\textwidth}
我门用一个两层网络来实现对这些离散点的拟合, 这两层网络的公式为:%
%
\begin{align*}
\bm{x}_2 &= \relu(\bm{W}_1x_1 + \bm{b}_1)\text{,} && \text{$\bm{W}_1\in\mathbb{R}^{10\times1}$, $\bm{b}_1\in\mathbb{R}^{10\times1}$,}\\
\hat{y} &= \bm{W}_2\bm{x}_2 + b_2\text{,}  && \text{$\bm{W}_2\in\mathbb{R}^{1\times10}$, $b_2\in\mathbb{R}^{1\times1}$.}
\end{align*}%
其中, 函数 $\relu(\,\cdot\,)$ 被称为线性整流函数(Rectified Linear Unit, ReLU), 图像如下所示.\vspace{-2pt}

\begin{center}
\hspace{-10pt}%
\begin{tikzpicture}
  \begin{axis}[width = 1.2\textwidth,
               height = 3.7cm,
               legend pos = north west,
  ]
    \addplot[mark = none, domain = -5:2] {(x < 0) * 0 + (x >= 0) * x};
    \legend{{$\mathrm{relu}(x) = \max{(0, x)}$}};

  \end{axis}
\end{tikzpicture}
\end{center}
\end{minipage}
\end{frame}

\begin{frame}[fragile]{\insertsection}{\insertsubsection}
\begin{pythoncode}[fontsize = \fontsize{8}{8}\selectfont, escapeinside=||]
import tensorflow as tf
import numpy as np

x1 = tf.placeholder(tf.float32, [None, 1])
w1 = tf.Variable(tf.random_normal([1, 10]))
b1 = tf.Variable(tf.ones([1, 10]))
# $\bm{x}_2 = \relu(\bm{W}_1x_1 + \bm{b}_1)$, $\bm{W}_1\in\mathbb{R}^{10\times1}$, $\bm{b}_1\in\mathbb{R}^{10\times1}$
x2 = tf.nn.relu(tf.matmul(x1, w1) + b1)
w2 = tf.Variable(tf.random_normal([10, 1]))
b2 = tf.Variable(tf.ones([1, 1]))
# $\hat{y} = \bm{W}_2\bm{x}_2 + b_2$, $\bm{W}_2\in\mathbb{R}^{1\times10}$, $b_2\in\mathbb{R}^{1\times1}$
predict_values = tf.matmul(x2, w2) + b2 # 预测值
actual_values = tf.placeholder(tf.float32, [None, 1]) # 真实值

# 目标函数 $\ell(\bm{W}_1, \bm{b}_1, \bm{W}_2, b_2) = \frac{1}{n}\sum_{i = 1}^{300}(y - \hat{y})^2$
loss = tf.reduce_mean(tf.reduce_sum(tf.square(actual_values - |\hspace{50pt}| predict_values), reduction_indices = [1]))
\end{pythoncode}

\vspace{-7.5cm}\hspace{10.5cm}\scalebox{0.53}{\begin{tikzpicture}[%
  tensor/.style = {rectangle, inner sep = 2pt, rounded corners = 0.4cm, draw, align = flush center, minimum height = 0.8cm, minimum width = 0.8cm},
  operator/.style = {tensor, fill = black, draw = none, text = white, inner sep = 2pt},
  line width = 1pt,
  y = 1.1cm
]
\node[tensor, circle]   (x1) at (9, 0) {$\bm{x}_1$};
\node[tensor, circle]   (w1) at (7, 0) {$\bm{W}_1$};
\node[operator]         (t1) at (8, 1) {\mintinline{text}{tf.matmul}};
\node[tensor, circle]   (b1) at (6, 1) {$\bm{b}_1$};
\node[operator]         (p1) at (7, 2) {\mintinline{text}{tf.add}};
\node[operator]         (r1) at (7, 3) {\mintinline{text}{tf.nn.relu}};
\node[tensor, circle]   (w2) at (5, 3) {$\bm{W}_2$};
\node[operator]         (t2) at (6, 4) {\mintinline{text}{tf.matmul}};
\node[tensor, circle]   (b2) at (4, 4) {$b_2$};
\node[operator]         (p2) at (5, 5) {\mintinline{text}{tf.add}};
\node[tensor]           (pv) at (5, 6) {\mintinline{text}{prediction_values}};
\node[tensor]           (av) at (8, 6) {\mintinline{text}{actual_values}};
\node[operator]         (sb) at (6.5, 7.5) {\mintinline{text}{tf.sub}};
\node[operator]         (sq) at (6.5, 8.5) {\mintinline{text}{tf.square}};
\node[operator]         (rs) at (6.5, 9.5) {\mintinline{text}{tf.reduce_sum}};
\node[operator]         (rm) at (6.5, 10.5) {\mintinline{text}{tf.reduce_mean}};
\node[tensor]           (ls) at (6.5, 11.5) {\mintinline{text}{loss}};

\draw[->] (x1) -- (t1);
\draw[->] (w1) -- (t1);
\draw[->] (b1) -- (p1);
\draw[->] (t1) -- (p1);
\draw[->] (p1) -- (r1);
\draw[->] (w2) -- (t2);
\draw[->] (r1) -- (t2);
\draw[->] (b2) -- (p2);
\draw[->] (t2) -- (p2);
\draw[->] (p2) -- (pv);
\draw[->] (pv) -- (sb);
\draw[->] (av) -- (sb);
\draw[->] (sb) -- (sq);
\draw[->] (sq) -- (rs);
\draw[->] (rs) -- (rm);
\draw[->] (rm) -- (ls);
\end{tikzpicture}}
\end{frame}

\begin{frame}[fragile]{\insertsection}{\insertsubsection}
\begin{pythoncode}[fontsize = \fontsize{8}{8}\selectfont]
def load_training_data(path): # 读取训练数据
    with open(path, 'r', encoding = 'utf8') as file:
        next(file)
        lines = [[float(item) for item in line.split(',')] for line in file]
    return (np.array([line[i] for line in lines])[:, None] for i in [0, 1])

with tf.Session() as session: # 初始化 Session
  session.run(tf.global_variables_initializer()) # 初始化变量

  (input_data, output_data) = load_training_data('./training_data.dat') # 加载训练数据
  train_step = tf.train.GradientDescentOptimizer(0.1).minimize(loss) # 设置目标函数

  for i in range(1000): # 训练 $1000$ 次
    session.run(train_step, feed_dict = {x1:input_data, actual_values:output_data})
    if i % 50 == 0:
      print(session.run(loss, feed_dict = {x1:input_data, actual_values:output_data}))
\end{pythoncode}
\end{frame}

\begin{frame}{\insertsection}{\insertsubsection}
训练 $1000$ 次后, 目标函数 $\ell(\bm{W}_1, \bm{b}_1, \bm{W}_2, b_2) = 2.48\times 10^{-3}$, 此时拟合函数的图像如图所示.\vspace{10pt}

\begin{tikzpicture}
\begin{axis}[width = 1\textwidth,
             height = 6.4cm,
             xlabel = $x$,
             ylabel = $y$,
             title = 训练数据与拟合曲线,
             legend pos = south east]
\addplot[only marks] table[col sep = comma] {./data/fitting/training_data.dat};
\addplot[mark = none,white, line width = 1pt, double = red, double distance = 2pt, on layer = foreground] table[col sep = comma, x = x, y = y] {./data/fitting/fitting_line.dat};
\legend{训练数据, 拟合曲线}
\end{axis}
\end{tikzpicture}
\end{frame}

% \begin{frame}[fragile]{\insertsection}{\insertsubsection}
% 我门用一个两层网络来实现对这些离散点的拟合, 这两层网络的公式为:%
% %
% \begin{flalign*}
%   && x_2 &= \relu(W_1x_1 + b_1)\text{,} && \text{$W_1\in\mathbb{R}^{10\times1}$, $b_1\in\mathbb{R}^{10\times1}$,} &&\\
%   && \hat{y} &= W_2x_2 + b_2\text{,}                  && \text{$W_2\in\mathbb{R}^{1\times10}$, $b_2\in\mathbb{R}^{1\times1}$.} &&
% \end{flalign*}

% \begin{minipage}[m]{0.38\textwidth}
% \begin{tikzpicture}
%   \begin{axis}[width = 1.2\textwidth,
%                height = 5.1cm,
%                legend pos = north west,
%   ]
%     \addplot[mark = none, domain = -2:5] {(x < 0) * 0 + (x >= 0) * x};
%     \legend{{$\mathrm{relu}(x) = \max{(0, x)}$}};

%   \end{axis}
% \end{tikzpicture}
% \end{minipage}%
% \hfill%
% \begin{minipage}[m]{0.60\textwidth}
% \begin{pythoncode}
% x1 = tf.placeholder(tf.float32, [None, 1])
% w1 = tf.Variable(tf.random_normal([1, 10]))
% b1 = tf.Variable(tf.ones([1, 10]))

% x2 = tf.nn.relu(tf.matmul(x1, w1) + b1)
% w2 = tf.Variable(tf.random_normal([10, 1]))
% b2 = tf.Variable(tf.ones([1, 1]))
% y = tf.matmul(x2, w2) + b2
% \end{pythoncode}
% \end{minipage}
% \end{frame}

% \subsection{训练网络}
% \begin{frame}[fragile]{\insertsection}{\insertsubsection}
% 我们用统计学中的均方误差作为代价函数来训练网络, 如下所示.%
% %
% \[
%   \ell(W_1, b_1, W_2, b_2) = \frac{1}{n}\sum_{i = 1}^n (y - \hat{y})^2\text{.}
% \]%
% %
% 根据上式可以写出 TensorFlow 中的训练目标.
% \begin{pythoncode}
% actual_values = tf.placeholder(tf.float32, [None, 1])
% loss = tf.reduce_mean(tf.reduce_sum(tf.square(actual_values - y), reduction_indices = [1]))

% train_step = tf.train.GradientDescentOptimizer(0.1).minimize(loss)
% \end{pythoncode}

% 我们可以看出, \inlinepython{tf.placeholder} 一般对应着训练数据, 因为训练数据只能在训练的时候给出, 在定义计算图的时候不知道具体的值, \inlinepython{tf.placeholder} 起占位的作用.
% \end{frame}

% \begin{frame}[fragile]{\insertsection}{\insertsubsection}
% 训练网络的过程, 就是求使代价函数 $\ell(W_1, b_1, W_2, b_2)$ 达到最小值时的参数集合 $W_1$, $b_1$, $W_2$, $b_2$. 而且需要启动一个会话 \inlinepython{tf.Session} 来完成运算图的计算.

% \begin{pythoncode}
% session = tf.Session()
% session.run(tf.initialize_all_variables())

% for step in range(1000):
%     feed_dict = {x1:input_data, actual_values:output_data}
%     session.run(train_step, feed_dict = feed_dict)

%     if step % 50 == 0:
%         print(session.run(loss, feed_dict = feed_dict))

% session.close()
% \end{pythoncode}
% \end{frame}

% \subsection{测试网络}

% \subsection{}

\insertthanks

\end{document}
